\documentclass{beamer}
\usepackage{amsmath}
\usepackage{amssymb}

% Choose a theme
\usetheme{Madrid}
\usecolortheme{default}
\usefonttheme{serif}

% --- CUSTOM FOOTER ---
% Redefine the footline to remove author/institute and keep title/page number
\makeatletter
\setbeamertemplate{footline}
{
  \leavevmode%
  \hbox{%
  \begin{beamercolorbox}[wd=.5\paperwidth,ht=2.25ex,dp=1ex,center]{title in head/foot}%
    \usebeamerfont{title in head/foot}\insertshorttitle
  \end{beamercolorbox}%
  \begin{beamercolorbox}[wd=.5\paperwidth,ht=2.25ex,dp=1ex,right]{date in head/foot}%
    \usebeamerfont{date in head/foot}\insertframenumber{} / \inserttotalframenumber\hspace*{2ex}
  \end{beamercolorbox}}%
  \vskip0pt%
}
\makeatother
% --- END CUSTOM FOOTER ---

\title[Lecture 16: Ray theory]{Lecture 16: Ray theory}
\author{David Al-Attar}
\institute{Michaelmas term 2024}
\date{}

\begin{document}

% --- INTRODUCTION ---

\begin{frame}
    \titlepage
\end{frame}

\begin{frame}{Outline and Motivation}
    In the last lecture, we studied plane waves in homogeneous media. What happens when the material properties vary with position?
    \begin{itemize}
        \item We will develop \textbf{ray theory}, an approximation for when the material changes slowly compared to the wavelength.
        \pause
        \item The key idea is that waves behave \textbf{locally} like plane waves, but their wavefronts curve and amplitudes change as they propagate.
        \pause
        \item This is a "high-frequency" approximation, analogous to \textbf{geometric optics} in electromagnetism.
    \end{itemize}
\end{frame}

% --- REGIMES ---

\begin{frame}{Wave Propagation Regimes}
    The behaviour of a wave interacting with a heterogeneous region depends on the length scales. Let $\lambda$ be the wavelength and $L$ be the scale of heterogeneity.
    \pause
    \begin{columns}
    \column{0.33\textwidth}
        \begin{alertblock}{Scattering}
            $L/\lambda \approx 1$
            \pause
            \centering
            The wave "sees" the object and scatters off it.
        \end{alertblock}

    \column{0.33\textwidth}
        \begin{alertblock}{Ray Theory}
            $L/\lambda \gg 1$
            \pause
            \centering
            The wave adapts smoothly to the slow changes. This is our focus.
        \end{alertblock}

    \column{0.33\textwidth}
        \begin{alertblock}{Averaging}
            $L/\lambda \ll 1$
            \pause
            \centering
            The wave doesn't resolve the fine details and "sees" an averaged medium.
        \end{alertblock}
    \end{columns}
\end{frame}

% --- RAY SERIES ANSATZ ---

\begin{frame}{Ray Theory Ansatz (1/2): Intuition}
    We assume the solution in the heterogeneous region takes a form similar to a plane wave, but with position-dependent terms.
    $$ u_{i}(\mathbf{x},t)=f[t-T(\mathbf{x})]a_{i}(\mathbf{x}) $$
    \pause
    \begin{block}{Travel Time $T(\mathbf{x})$}
        The phase is now a non-linear function of position. The surfaces $T(\mathbf{x}) = \text{constant}$ are the \textbf{wavefronts}. The gradient of the travel time defines the local \textbf{slowness vector}:
        $$ p_i(\mathbf{x}) = \frac{\partial T}{\partial x_i} $$
    \end{block}
    \pause
    \begin{block}{Amplitude $a_i(\mathbf{x})$}
        The amplitude and polarisation vector can now vary in space, accounting for focusing/defocusing and changes in particle motion.
    \end{block}
\end{frame}

\begin{frame}{Ray Theory Ansatz (2/2): The Ray Series}
    To be more rigorous, we seek a solution as a series expansion in inverse powers of frequency, $\omega$. In the frequency domain, this is:
    $$ u_{i}(\mathbf{x},\omega)=\left(\sum_{n=0}^{\infty}\frac{a_{i}^{n}(\mathbf{x})}{(i\omega)^{n}}\right)e^{-i\omega T(\mathbf{x})} $$
    \pause
    \begin{itemize}
        \item This is the \textbf{elastodynamic ray series}.
        \pause
        \item The $n=0$ term corresponds to our initial guess. The higher-order terms are progressively smoother corrections.
        \pause
        \item This is an \textbf{asymptotic expansion}, which is accurate for high frequencies ($\lambda \ll L$) but may not converge in the usual sense.
    \end{itemize}
\end{frame}

\begin{frame}{Hierarchy of Equations}
    Substituting the ray series into the equations of motion and collecting terms with the same power of $i\omega$ gives a hierarchy of equations.
    \pause
    \begin{block}{Highest Order ($ (i\omega)^2 $): The Eikonal Equation}
        This equation determines the travel time, $T(\mathbf{x})$.
        $$ [A_{ijkl}p_{j}p_{l}-\rho\delta_{ik}]a_{k}^{0}=0 $$
        where $p_j = \partial T / \partial x_j$.
    \end{block}
    \pause
    \begin{block}{Next Order ($ (i\omega)^1 $): The Transport Equation}
        This equation determines the amplitude of the leading-order term, $a_k^0$.
        $$ [A_{ijkl}p_{j}p_{l}-\rho\delta_{ik}]a_{k}^{1}=\text{Terms involving } a_k^0 $$
    \end{block}
    \pause
    We will focus on solving the first and most important of these: the Eikonal equation.
\end{frame}

% --- EIKONAL EQUATION ---

\begin{frame}{The Eikonal Equation}
    Let's look at the highest-order equation from the ray series:
    $$ [A_{ijkl}p_{j}p_{l}-\rho\delta_{ik}]a_{k}^{0}=0 $$
    \pause
    \begin{itemize}
        \item This looks exactly like the Christoffel equation for plane waves, but now the material properties ($\rho$, $A_{ijkl}$) and the slowness vector ($\mathbf{p}$) depend on position $\mathbf{x}$.
        \pause
        \item For a non-trivial solution for the polarisation $a_k^0$, the determinant of the matrix must be zero.
    \end{itemize}
    \pause
    \begin{alertblock}{The Eikonal Equation}
        $$ \det[\mathbf{\Gamma}(\mathbf{x}, \mathbf{p}(\mathbf{x}))-\mathbf{I}]=0 $$
        This is a first-order, non-linear partial differential equation for the travel time $T(\mathbf{x})$, since $\mathbf{p} = \nabla T$.
    \end{alertblock}
\end{frame}

\begin{frame}{Solving the Eikonal Equation}
    \begin{exampleblock}{Geometric Interpretation}
        The eikonal equation means that at every point $\mathbf{x}$, the local slowness vector $\mathbf{p}(\mathbf{x})$ must lie on the local slowness surface.
    \end{exampleblock}
    \pause
    \begin{itemize}
        \item If we are tracking a single wave type (e.g., a quasi-P wave), the eikonal equation simplifies.
        \pause
        \item Let $c_k(\mathbf{x}, \hat{\mathbf{p}})$ be the phase speed for the $k$-th wave type in direction $\hat{\mathbf{p}}$ at point $\mathbf{x}$. The equation becomes:
        $$ ||\mathbf{p}(\mathbf{x})||^2 c_k[\mathbf{x}, \hat{\mathbf{p}}(\mathbf{x})]^2 = 1 $$
    \end{itemize}
    \pause
    To solve this PDE, we use a powerful technique called the \textbf{method of characteristics}.
\end{frame}

% --- METHOD OF CHARACTERISTICS ---

\begin{frame}{Method of Characteristics (1/2): Rays}
    The method of characteristics turns the PDE for $T(\mathbf{x})$ into a system of ODEs that trace out rays.
    \begin{itemize}
        \item First, we define the \textbf{Ray Hamiltonian} for the $k$-th wave type:
        $$ H_{k}(\mathbf{x},\mathbf{p}) = \frac{1}{2}||\mathbf{p}||^{2}c_{k}[\mathbf{x},\hat{\mathbf{p}}]^{2} $$
        The eikonal equation is then simply $H_k = 1/2$.
        \pause
        \item The path of a ray, $\mathbf{x}(\sigma)$, and its evolving slowness vector, $\mathbf{p}(\sigma)$, are found by integrating \textbf{Hamilton's equations}:
    \end{itemize}
    \begin{alertblock}{Hamilton's Ray Tracing Equations}
        $$ \frac{d x_{i}}{d\sigma}=\frac{\partial H_{k}}{\partial p_{i}} \quad \text{and} \quad \frac{d p_{i}}{d\sigma}=-\frac{\partial H_{k}}{\partial x_{i}} $$
    \end{alertblock}
\end{frame}

\begin{frame}{Method of Characteristics (2/2): Travel Time}
    \begin{block}{Physical Meaning}
        \begin{itemize}
            \item The curves $\mathbf{x}(\sigma)$ traced by Hamilton's equations are the \textbf{rays} — the paths that wave energy follows.
            \pause
            \item The integration parameter $\sigma$ has a direct physical meaning: it is the \textbf{travel time} along the ray.
            $$ T[\mathbf{x}(\sigma)] = \sigma $$
        \end{itemize}
    \end{block}
    \pause
    \begin{block}{Caustics}
        \begin{itemize}
            \item A \textbf{caustic} is a location where rays cross.
            \pause
            \item At a caustic, the travel time function is multi-valued, and the ray theory amplitude $a_k^0$ becomes infinite.
            \pause
            \item This is a breakdown of the approximation. In reality, amplitudes at caustics are very large but finite.
        \end{itemize}
    \end{block}
\end{frame}

% --- SUMMARY ---

\begin{frame}{Summary: What You Need to Know}
    \begin{itemize}
        \item Ray theory is a high-frequency approximation ($L \gg \lambda$) used for smoothly varying media.
        \pause
        \item The \textbf{ray series} is an asymptotic expansion for the wavefield. Substituting it into the wave equation gives a hierarchy of equations.
        \pause
        \item The highest-order equation is the \textbf{eikonal equation}, a PDE for the travel time $T(\mathbf{x})$.
        \pause
        \item The eikonal equation is solved by the \textbf{method of characteristics}, which involves integrating \textbf{Hamilton's equations} to trace \textbf{rays}.
        \pause
        \item The parameter along a ray is the travel time. A \textbf{caustic} is formed where rays cross, and is a point where simple ray theory breaks down.
    \end{itemize}
\end{frame}

\end{document}
